\documentclass[10pt,a4paper]{amsart}
\usepackage[margin=19mm]{geometry}
\usepackage{amsmath,amssymb,mathtools}
\usepackage[scr=rsfs,cal=euler]{mathalfa}
\usepackage{xfrac}
\usepackage[unicode,hidelinks,bookmarks=false]{hyperref}
\newcommand{\ie}{i.e.\ }
\newcommand{\cf}{cf.\ }
\newcommand{\resp}{respectively\ }
\newcommand{\Subsection}{\subsection}
\providecommand{\nobreakdash}{\nobreak\hbox{-}}
\setlength{\emergencystretch}{2em}
\multlinegap0pt
\usepackage{amssymb}
\usepackage{algorithm}\usepackage{algpseudocode}
\usepackage{enumerate}
\usepackage{float}
\newtheorem{prop}{Proposition}
\newtheorem{thm}{Theorem}
\title{Hamiltonian Interface Dynamics for Reduced-Order Optimization of Incompressible Mixing}
\author{Ziqian Li, Enrique Zuazua}
\date{Source: arXiv:2605.04688v1 (CC BY 4.0)}
\begin{document}
\maketitle

\noindent\emph{Source and licence.} Hamiltonian Interface Dynamics for Reduced-Order Optimization of Incompressible Mixing. Version arXiv:2605.04688v1 (CC BY 4.0), distributed by arXiv under the Creative Commons Attribution 4.0 International licence, http://creativecommons.org/licenses/by/4.0/, as recorded in the arXiv licence metadata for this exact version and shown on the abstract page https://arxiv.org/abs/2605.04688v1. Full licence notice: \texttt{licenses/arxiv-2605.04688-CC-BY-4.0.txt}.\par\smallskip

\noindent\textit{Editorial scope.} Counted unit: the article's prop prop:adjoint (source0.tex lines 556--568). The article's complete original proof of it is delivered with it (source0.tex lines 570--614, one \texttt{proof} environment of 45 lines). No mathematics is written, repaired or re-derived: the statement and the proof are the article's own lines, in the article's own notation, and no retained line is substituted or re-flowed.  Retained uncounted as the source-local context the proof needs: lines 224--237; lines 265--267; lines 269--271; lines 273--275; lines 279--285; lines 287--289; lines 298--300; lines 303--309; lines 313--315; lines 353--364; lines 427--429; lines 530--537.  Nothing was omitted from any retained span, so the package carries no omission record.  No retained line contains a citation, so the wrapper carries no bibliography and no bibliography entry.  No retained line contains a figure environment, an image-inclusion command or drawing code, so no image is delivered and the package declares no asset dependency.  Editorial formatting only, in addition: an AMS \texttt{amsart} preamble that restates the article's own declarations and macros used by the retained lines, namely the environments \texttt{prop}, \texttt{thm} on separate counters, together with the standard packages those lines need.  The retained environments print in this document as Proposition 1, Theorem 1 in order of appearance, whereas the article prints the same items with the numbering of its own full document.  The complete native source of the article as distributed in the arXiv e-print for this exact version is bundled unchanged as \texttt{sources/199883/source0.tex}.
\begin{prop}[Interface dynamics]\label{prop:interface_ode}
The pair $(X,a)$ satisfies the coupled ODE system
\begin{subequations}\label{eq:state}
\begin{align}
\partial_t X(t,s) &= J\nabla_x \psi\bigl(t,X(t,s)\bigr),\label{eq:state_X}\\
\partial_t a(t,s) &= J\nabla_x^2 \psi\bigl(t,X(t,s)\bigr)\,a(t,s),\label{eq:state_a}\\
X(0,s) &= \gamma_0(s),\quad a(0,s) = \gamma_0'(s).
\end{align}
\end{subequations}
The length of the advected interface at time $t$ is
\begin{equation}\label{eq:length}
\mathcal{L}(\psi;t):=\int_0^1 |a(t,s)|\,ds.
\end{equation}
\end{prop}

\begin{equation}\label{eq:basis_assumption}
h_k\in C^\infty(\overline\Omega)\cap H_0^1(\Omega),\qquad k=1,\dots,N,
\end{equation}

\begin{equation}\label{eq:smooth_space}
\mathcal H_N:=\mathrm{span}\{h_1,\dots,h_N\}\subset C^\infty(\overline\Omega)\cap H_0^1(\Omega).
\end{equation}

\begin{equation}\label{eq:basis_constants}
B^{(j)}:=\max_{1\le k\le N}\|h_k\|_{W^{j,\infty}(\Omega)},\qquad j=0,1,2,3,
\end{equation}

\begin{equation}\label{eq:Uad}
\mathcal U_{\mathrm{ad}}^{R_u}
:=
\Bigl\{
\psi(t,x)=\sum_{k=1}^N u_k(t)\,h_k(x)\;;\;u\in L^2(0,T;\mathbb R^N),\ \|u\|_{L^2(0,T;\mathbb R^N)}\le R_u
\Bigr\}.
\end{equation}

\begin{equation}\label{eq:pt_bound}
\|\psi(t,\cdot)\|_{W^{j,\infty}(\Omega)}\le\sqrt N\,B^{(j)}\,|u(t)|\qquad\text{for a.e. }t\in(0,T),
\end{equation}

\begin{equation}\label{eq:gram_identity}
\|\psi\|_{L^2(\Omega)}^2=u^\top M u.
\end{equation}

\begin{equation}\label{eq:cost}
\mathcal{J}_{\lambda,\varepsilon}(\psi)
:=
-\int_0^1 \ell_\varepsilon\bigl(a(T,s)\bigr)\,ds
+\frac{\lambda}{2}\int_0^T \|\psi(t,\cdot)\|_{L^2(\Omega)}^2\,dt,
\qquad\psi\in\mathcal U_{\mathrm{ad}}^{R_u},
\end{equation}

\begin{equation}\label{eq:ocp}
\min_{\psi\in\mathcal{U}_{\mathrm{ad}}^{R_u}}\mathcal{J}_{\lambda,\varepsilon}(\psi).
\end{equation}

\begin{thm}[Existence of a minimizer]\label{prop:H3}
Let $\mathcal H_N$, $B^{(j)}$, $\mathcal U_{\mathrm{ad}}^{R_u}$, and $\mathcal J_{\lambda,\varepsilon}$ be given by \eqref{eq:smooth_space}, \eqref{eq:basis_constants}, \eqref{eq:Uad}, and \eqref{eq:cost}, respectively. Assume that the basis $h_k$ satisfies \eqref{eq:basis_assumption} and that $\gamma_0\in W^{2,\infty}([0,1];\overline\Omega)$. Then the following properties hold.
\begin{enumerate}
\item[\textup{(i)}]\emph{(State regularity and well-posedness.)} For every $u\in L^2(0,T;\mathbb R^N)$ with associated stream function $\psi=\sum_{k=1}^N u_k h_k$,
\begin{equation}\label{eq:v_reg}
\psi\in L^2(0,T;W^{3,\infty}(\overline\Omega)),\qquad
\mathbf v=\nabla^\perp\psi\in L^2(0,T;W^{2,\infty}(\Omega;\mathbb R^2));
\end{equation}
the flow $\Phi_\psi$ maps $\overline\Omega$ into itself, and the state system  \eqref{eq:state} admits a unique solution $(X,a)\in C([0,T]\times[0,1];\overline\Omega\times\mathbb R^2)$. Consequently, the cost functional $\mathcal{J}_{\lambda,\varepsilon}$ is well defined and depends continuously on $u$, i.e., on  $\psi$.
\item[\textup{(ii)}]\emph{(Existence of a minimizer.)} The minimization problem \eqref{eq:ocp} admits at least one minimizer $\psi^*\in\mathcal U_{\mathrm{ad}}^{R_u}$.
\end{enumerate}
\end{thm}

\begin{equation}\label{eq:Ma}
|a(t,s)|\le M_a:=\|\gamma_0'\|_{L^\infty([0,1])}\,\exp\,\bigl(\sqrt N\,B^{(2)}\,\sqrt T\,R_u\bigr)
\end{equation}

\begin{align}
-\partial_t p(t,s)
&=J\nabla_x^2\psi(t,X)^\top p(t,s)+\bigl(J\nabla_x^3\psi(t,X)[\,\cdot\,,a(t,s)]\bigr)^{\!\top}\!q(t,s),\label{eq:adjoint_p}\\
-\partial_t q(t,s)
&=J\nabla_x^2\psi(t,X)^\top q(t,s),\label{eq:adjoint_q}\\
p(T,s)&=0,\qquad
q(T,s)=-\nabla\ell_\varepsilon(a(T,s))=-\frac{a(T,s)}{\sqrt{|a(T,s)|^2+\varepsilon^2}},\label{eq:adjoint_tc}
\end{align}

\begin{prop}[Reduced gradient via the adjoint system]\label{prop:adjoint}
Under the hypotheses of Theorem~\ref{prop:H3} and set $j_{\lambda,\varepsilon}(u):=\mathcal J_{\lambda,\varepsilon}(\psi)$. Then $j_{\lambda,\varepsilon}$ is Fr\'echet differentiable on $L^2(0,T;\mathbb R^N)$, and for every $\delta u\in L^2(0,T;\mathbb R^N)$ with $\eta:=\sum_{k=1}^N\delta u_k h_k$,
\begin{align}\label{eq:DJ}
Dj_{\lambda,\varepsilon}(u)[\delta u]
=\lambda\!\int_0^T\!\!(Mu(t))\cdot\delta u(t)\,dt
+\!\int_0^T\!\!\int_0^1\!\Bigl[p\cdot J\nabla_x\eta(t,X)+q\cdot J\nabla_x^2\eta(t,X)\,a\Bigr]ds\,dt,
\end{align}
or, componentwise, with $k=1,\dots,N$,
\begin{equation}\label{eq:grad_uk}
\frac{\partial j_{\lambda,\varepsilon}}{\partial u_k}(t)
=\lambda\,(Mu(t))_k+\int_0^1\!\Bigl[p(t,s)\cdot J\nabla h_k(X)+q(t,s)\cdot J\nabla_x^2 h_k(X)\,a(t,s)\Bigr]ds.
\end{equation}
\end{prop}

\begin{proof}
\textit{Step 1. Linearized state.}
Since each $h_k\in C^\infty(\overline\Omega)$, the right-hand side of \eqref{eq:state} is $C^\infty$ in $(X,a)$ and affine in $u$. Gr\"onwall's inequality applied to the difference of two state trajectories, with the $L^2$-integrable Lipschitz constants furnished by the pointwise bound \eqref{eq:pt_bound} and \eqref{eq:Ma}, yields
\[
\|(X_{u+\delta u}-X_u,\,a_{u+\delta u}-a_u)\|_{C([0,T]\times[0,1])}\le C\,\|\delta u\|_{L^2(0,T;\mathbb R^N)},
\]
and a second Gr\"onwall pass on the quadratic remainder gives an $O(\|\delta u\|_{L^2}^2)$ bound. Hence $u\mapsto(X,a)$ is Fr\'echet differentiable from $L^2(0,T;\mathbb R^N)$ into $C([0,T]\times[0,1];\mathbb R^2\times\mathbb R^2)$, and its derivative $(\delta X,\delta a):=(DX[\delta u],Da[\delta u])$ solves
\begin{subequations}
\begin{align}
\partial_t\delta X &= J\nabla_x^2\psi(t,X)\,\delta X+\sum_{k=1}^N\delta u_k(t)\,J\nabla h_k(X),\label{eq:linearized_X}\\
\partial_t\delta a &= J\nabla_x^2\psi(t,X)\,\delta a+J\nabla_x^3\psi(t,X)[\delta X,a]+\sum_{k=1}^N\delta u_k(t)\,J\nabla_x^2 h_k(X)\,a,\label{eq:linearized_a}\\
\delta X(0,s)&=0,\qquad \delta a(0,s)=0.\label{eq:linearized_ic}
\end{align}
\end{subequations}

\vskip 5pt
\noindent\textit{Step 2. Lagrangian identity.}
Pair \eqref{eq:linearized_X} with $p$ and \eqref{eq:linearized_a} with $q$ in $L^2((0,T)\times(0,1))$ and integrate by parts in $t$; the boundary conditions \eqref{eq:linearized_ic} and $p(T,\cdot)=0$ of \eqref{eq:adjoint_tc} give
\begin{equation}\label{eq:IBP}
\begin{aligned}
\int_0^T\!\!\int_0^1 p\cdot\partial_t\delta X\,ds\,dt &= -\int_0^T\!\!\int_0^1 \partial_t p\cdot\delta X\,ds\,dt,\\
\int_0^T\!\!\int_0^1 q\cdot\partial_t\delta a\,ds\,dt &= \int_0^1 q(T,s)\cdot\delta a(T,s)\,ds-\int_0^T\!\!\int_0^1 \partial_t q\cdot\delta a\,ds\,dt.
\end{aligned}
\end{equation}
Substituting \eqref{eq:linearized_X}--\eqref{eq:linearized_a} on the left and \eqref{eq:adjoint_p}--\eqref{eq:adjoint_q} on the right of \eqref{eq:IBP}, the transpose identities
\[
p\cdot J\nabla_x^2\psi\,\delta X=(J\nabla_x^2\psi)^\top p\cdot\delta X,
\qquad
q\cdot J\nabla_x^3\psi[\delta X,a]=\bigl(J\nabla_x^3\psi[\,\cdot\,,a]\bigr)^\top q\cdot\delta X,
\]
and the analogous identity for $q\cdot J\nabla_x^2\psi\,\delta a$, pair up the bulk contributions and cancel them. What remains is
\begin{align}\label{eq:Lagrangian_id}
\int_0^1 q(T,s)\cdot\delta a(T,s)\,ds
=\int_0^T\!\!\int_0^1\sum_{k=1}^N\delta u_k(t)\Bigl[p\cdot J\nabla h_k(X)+q\cdot J\nabla_x^2 h_k(X)\,a\Bigr]ds\,dt.
\end{align}

\vskip 5pt
\noindent\textit{Step 3. Differentiating the cost.}
By \eqref{eq:cost} and \eqref{eq:gram_identity}, $$\mathcal J_{\lambda,\varepsilon}(\psi)=\frac\lambda2\int_0^T u^\top Mu\,dt-\int_0^1\ell_\varepsilon(a(T,s))\,ds.$$ Differentiating in $\delta u$ and using $q(T,s)=-\nabla\ell_\varepsilon(a(T,s))$,
\[
Dj_{\lambda,\varepsilon}(u)[\delta u]
=\lambda\!\int_0^T\!\!(Mu)\cdot\delta u\,dt+\int_0^1 q(T,s)\cdot\delta a(T,s)\,ds,
\]
which combined with \eqref{eq:Lagrangian_id} yields \eqref{eq:DJ}. Taking $\delta u_\ell(t)=\delta_{k\ell}\varphi(t)$ for arbitrary $\varphi\in L^2(0,T)$ specializes \eqref{eq:DJ} to \eqref{eq:grad_uk}.
\end{proof}

\end{document}
